% Consolidated 07-hydrodynamics from 6915d31494b4647dd12b50a59cddbc000c3484fc; original TeX fragments preserved.

% === BEGIN 07-hydrodynamics-315-324.tex blob cc9c9e4fa0ffd378e3fc3aa288b48e691f93bace ===
% User-provided Irodov PDF pp.59-61, original printed pp.50-52.
% Original section 1.7: hydrodynamics, Figs. 1.80-1.86 checked.
\exproblem{315}{51}
理想不可压缩流体沿一条置于水平面内、
横截面积恒定但呈直角弯曲的扁平管稳定流动
（原书图1.80，俯视图）。
图中弯道外侧位置为点1、内侧位置为点2。
问两点压强和速率是否相同，并描述流线疏密。
\exsolution
在弯管内，流体质元须具有指向弯道曲率中心的
法向加速度 $v^2/R_c$。沿弯道横向法线，
稳态欧拉方程近似写作
\[
\frac{\partial p}{\partial n}
=\rho\frac{v^2}{R_c},
\]
其中 $n$ 指向弯道外缘。
因此压强从内缘向外缘升高，
即 $\boxed{p_1>p_2}$。
在图示的扁平弯道中，流线向内弯侧挤密，
给定相邻流管的流量不变，流线越密处局部速度越大，
因而图示 $\boxed{v_1<v_2}$。
流线大体顺着弯道从入口转向出口，
内缘间距较小，外缘较疏。
\excheck 这里的两点位于\emph{不同流线}，
不可不加区分地拿“同一条流线上”的伯努利常数
直接比较 $p_1$ 与 $p_2$。
这一判断与原图中的流线分布一致。

\exproblem{316}{51}
一根水平变截面管道上在截面 $S_1,S_2$
分别安装竖直测压管（原书图1.81），两管水柱高度差
为 $\Delta h$。假设水为理想不可压缩流体、流动稳态。
求管中体积流量 $Q$。
\exsolution
连续性方程给出
$v_1=Q/S_1$、$v_2=Q/S_2$。
两截面等高，沿流线应用伯努利方程：
\[
p_1+\frac12\rho\frac{Q^2}{S_1^2}
=p_2+\frac12\rho\frac{Q^2}{S_2^2}.
\]
测压管显示
$|p_1-p_2|=\rho g\Delta h$。故
\[
g\Delta h=\frac{Q^2}{2}
\left|\frac1{S_1^2}-\frac1{S_2^2}\right|.
\]
解出
\[
\boxed{Q=S_1S_2\sqrt{
\frac{2g\Delta h}{|S_2^2-S_1^2|}}}.
\]
\excheck 两截面必须不等，否则压强差对应流量无法由
该方程唯一求出；若已知
$S_2>S_1$，可直接去掉绝对值为
$\sqrt{2g\Delta h/(S_2^2-S_1^2)}$。

\exproblem{317}{51}
气体在截面积为 $S$ 的管中稳定流动，
中心装有皮托管（原书图1.82）。
差压计液体密度 $\rho_0$，气体密度 $\rho$，
液柱高度差为 $\Delta h$。忽略气体黏性，
求体积流量 $Q$，并说明近似条件。
\exsolution
皮托管迎流孔测停滞压强，
静压孔测流动压强；同高测点的差压为
\[
p_{\rm stag}-p_{\rm static}=\frac12\rho v^2.
\]
U形差压计若两侧气体等密度，其两端压差为
$\Delta p=(\rho_0-\rho)g\Delta h$。
二式相等：
\[
\frac12\rho v^2=(\rho_0-\rho)g\Delta h,
\]
因此
\[
\boxed{Q=Sv=S\sqrt{
\frac{2g\Delta h(\rho_0-\rho)}{\rho}}}.
\]
通常 $\rho_0\gg\rho$，则
\[
\boxed{Q\simeq S\sqrt{\frac{2g\Delta h\rho_0}{\rho}}}.
\]
\excheck 以截面中心速度代替截面平均速度，
还须假定截面内流速分布近似均匀；
实际管流有速度剖面时仅由中心皮托测量不能自动
确定精确总流量。

\exproblem{318}{51}
一个横截面很大的容器中，下层为水、
上层为煤油。
水层厚度 $h_1=30\,\mathrm{cm}$，
煤油层厚度 $h_2=20\,\mathrm{cm}$；
容器底部有小孔通向大气。
忽略黏性，求流出水的速度。
\exsolution
记水和煤油密度分别为 $\rho_1,\rho_2$。
容器很宽，内部两层自由表面下降速度可忽略。
在容器底部小孔处，来自液柱的表压为
\[
p-p_0=\rho_1gh_1+\rho_2gh_2.
\]
在小孔出口，水的静压为大气压 $p_0$。
将液柱压能转化为水的动能：
\[
\frac12\rho_1v^2
=\rho_1gh_1+\rho_2gh_2.
\]
故
\[
\boxed{v=\sqrt{2g\left(
h_1+\frac{\rho_2}{\rho_1}h_2\right)}}.
\]
煤油相对密度取约 $0.80$，得到
$v\simeq\sqrt{2(9.81)(0.30+0.80\cdot0.20)}
=\boxed{3.0\,\mathrm{m/s}}$。
\excheck 不可把上方煤油层直接当成同厚度的水层；
它的贡献须乘密度比 $\rho_2/\rho_1$。

\exproblem{319}{51}
高为 $H=50\,\mathrm{cm}$ 的宽圆柱形水箱装满水，
放在桌面上（箱底即桌面高度）。
在侧壁距底部高度为 $y$ 的位置钻一小孔，
水流水平射出并落到桌面。
忽略黏性，求使落点距箱壁最远的孔高
$y_{\rm opt}$ 与最大水平射程。
\exsolution
自由液面到孔的水头为 $H-y$，
托里拆利公式
$v=\sqrt{2g(H-y)}$。
孔到桌面垂直距离为 $y$，
水平抛体运动的落地时间为
$t=\sqrt{2y/g}$。
射程
\[
L(y)=vt=2\sqrt{y(H-y)},\qquad 0<y<H.
\]
由
$y(H-y)=H^2/4-(y-H/2)^2$
知在 $y=H/2$ 取得全局最大值：
\[
\boxed{y_{\rm opt}=H/2=25\,\mathrm{cm}},\qquad
\boxed{L_{\max}=H=50\,\mathrm{cm}}.
\]
\excheck 孔高接近底部时落地时间太短，
接近水面时喷出速度太小，两端的射程均趋零。

\exproblem{320}{51}
按原书图1.83，弯曲细管迎着水平水流放置，
水流相对细管速率 $v=2.5\,\mathrm{m/s}$。
细管露出水面并在高于水面
$h_0=12\,\mathrm{cm}$ 的闭口端顶部开一个小喷孔。
求喷出的水能够达到喷孔上方的高度 $h$。
\exsolution
迎流管入口将流动动压
$\tfrac12\rho v^2$ 转成停滞压力；
等效水柱高度为
$H_{\rm dyn}=v^2/(2g)$。
管内水从水面上升 $h_0$
消耗重力势能，余下的水头在喷孔处形成喷出速度，
并转为喷孔上方上升高度 $h$。
伯努利方程给出
\[
h+h_0=\frac{v^2}{2g}.
\]
因此
\[
\boxed{h=\frac{v^2}{2g}-h_0
=\frac{(2.5)^2}{2(9.81)}-0.12
\simeq0.20\,\mathrm m=20\,\mathrm{cm}}.
\]
\excheck 必须有 $v^2/(2g)>h_0$
才能从顶端喷出；否则水面仅上升但不出孔。

\exproblem{321}{52}
宽水槽底部开半径为 $R_1$ 的圆孔，
孔上方有半径 $R_2>R_1$ 的封闭圆柱盘，
盘与槽底之间留极窄的水平径向间隙
（原书图1.84）。上部自由水面高于底部 $h$，
液体密度 $\rho$。
求圆柱盘下方间隙内距离圆心 $r$ 处的静压强，
其中 $R_1<r<R_2$。
\exsolution
设外部大气压为 $p_0$。
宽槽静止水体的总水头相当于压强
$p_0+\rho gh$。
在极窄而恒厚度的间隙中，
水沿径向向孔口流动；不可压缩连续性使
$r v_r(r)=$常数。
把小孔边缘 $r=R_1$ 处静压近似取为 $p_0$，
由储水液面的伯努利关系，
其间隙流速大小可取
$v_r(R_1)=\sqrt{2gh}$。
故
\[
v_r(r)=\frac{R_1}{r}\sqrt{2gh}.
\]
对间隙内任意点，
伯努利方程为
$p(r)+\tfrac12\rho v_r(r)^2=p_0+\rho gh$，
于是
\[
\boxed{p(r)=p_0+\rho gh
\left(1-\frac{R_1^2}{r^2}\right)},
\qquad R_1<r<R_2.
\]
\excheck 在 $r=R_1$ 边缘 $p=p_0$，
远离孔口时 $p$ 向静水压 $p_0+\rho gh$ 靠近；
该理想模型忽略极窄缝的黏性损失，
现实中的黏性阻力通常显著。

\exproblem{322}{52}
原书图1.85显示，水平圆筒内有体积 $V$ 的水，
活塞受\emph{恒定}外推力，
水经远小于活塞面积的小孔 $s$ 喷出。
假设无摩擦、无黏性，
要在时间 $t$ 内将水全部挤出，至少需要做多少功？
\exsolution
由水的不可压缩性和题定恒力理想化，
若流量恒定，则
$Q=V/t$，小孔平均喷射速率为
$v=Q/s=V/(st)$。
小孔射出的水动能累积为
\[
K_{\rm out}=
\frac12\rho V v^2
=\frac12\rho V\left(\frac{V}{st}\right)^2.
\]
忽略活塞及筒内水的末态动能、
大气压两侧抵消、没有机械损失，则外力最小功
\[
\boxed{W=\frac{\rho V^3}{2s^2t^2}}.
\]
也可以由所需压差
$\Delta p=\rho v^2/2$ 乘排出体积 $V$ 得同式。
\excheck 恒流量近似要求筒内入口截面远大于小孔截面；
若保留活塞及筒内水的非稳态加速过程，
还需要额外分析瞬态动能与力的时间规律。

\exproblem{323}{52}
一个底面积为 $S$、高度为 $h$ 的竖直圆筒装满水，
底部开面积为 $s\ll S$ 的小孔。
忽略流动黏性，求水全部流出的时间 $\tau$。
\exsolution
设瞬时水面高度为 $z(t)$，
小孔的流出速度由托里拆利关系为
$v(z)=\sqrt{2gz}$。
不可压缩连续性要求箱中体积减少率等于出流量：
\[
-S\frac{\mathrm dz}{\mathrm dt}
=s\sqrt{2gz}.
\]
分离变量积分：
\[
\int_h^0\frac{\mathrm dz}{\sqrt z}
=-\frac{s}{S}\sqrt{2g}\int_0^\tau\mathrm dt.
\]
左侧为 $-2\sqrt h$，
故
\[
\boxed{\tau=\frac{S}{s}\sqrt{\frac{2h}{g}}}.
\]
\excheck $s\ll S$ 是忽略水面下落速度
在伯努利方程中贡献的关键；
实际黏性和孔口收缩会延长排水时间。

\exproblem{324}{52}
长为 $\ell$ 的水平管 $AB$ 绕其一端 $A$
处固定竖直轴以角速度 $\omega$ 匀速旋转
（原书图1.86）。
靠轴端 $A$ 开口，远端 $B$ 封闭但有微小出流孔。
理想液体在管内形成从 $B$ 向 $A$
长度为 $h$ 的液柱，即靠内侧自由界面处于
$x=\ell-h$。求液体相对管道的出流速度。
\exsolution
在随管匀速转动的参考系中，
单位质量的离心势为 $-\omega^2x^2/2$。
液柱内端 $x_A=\ell-h$
与出流孔端 $x_B=\ell$ 都接大气压，
忽略黏性及前端自由界面速度，
沿径向相对流线应用旋转系伯努利积分：
\[
\frac12v^2
=\frac{\omega^2}{2}
\bigl[x_B^2-x_A^2\bigr]
=\frac{\omega^2}{2}\bigl[\ell^2-(\ell-h)^2\bigr].
\]
因此
\[
\boxed{v=\omega\sqrt{2\ell h-h^2}
=\omega h\sqrt{\frac{2\ell}{h}-1}},
\qquad 0<h\le\ell.
\]
\excheck $v$ 受管内液柱长度 $h$ 控制；
题目中的 $h$ 是\emph{水平径向液柱长度}，
不能误解为竖直静水高度。
% === END 07-hydrodynamics-315-324.tex ===

% === BEGIN 07-hydrodynamics-325-334.tex blob 858e460f1eead93526268874f67eb6d153c555bf ===
% Irodov user PDF, printed pp.52-54, PDF pp.61-63,
% hydrodynamics problems 325-334, Figs 1.87-1.89 inspected.
\exproblem{325}{52}
从理想流体的欧拉运动方程出发，
证明在不可压缩流体的定常流动中，
沿任一流线成立伯努利方程。
\exsolution
理想流体在重力场中的欧拉方程为
\[
\rho\frac{D\bm v}{Dt}=-\nabla p+\rho\bm g.
\]
定常流动时
$\partial\bm v/\partial t=0$，因此
$D\bm v/Dt=(\bm v\cdot\nabla)\bm v$。
令流线弧长为 $s$、单位切向矢量
$\bm t=\bm v/v$，
把上式与 $\bm t$ 作点积：
\[
\rho v\frac{\mathrm dv}{\mathrm ds}
=-\frac{\mathrm dp}{\mathrm ds}
-\rho g\frac{\mathrm dz}{\mathrm ds}.
\]
这里竖直坐标 $z$ 向上，故
$\bm g\cdot\bm t=-g\,\mathrm dz/\mathrm ds$。
在 $\rho$ 为常数的不可压缩近似下，
\[
\frac{\mathrm d}{\mathrm ds}
\left(\frac12\rho v^2+p+\rho gz\right)=0.
\]
对同一条流线积分：
\[
\boxed{\frac12\rho v^2+p+\rho gz=C_{\rm streamline}}.
\]
\excheck 每条流线的常数一般可能不同；
只有特定条件（如无旋势流）才能进一步说明不同
流线拥有同一个伯努利常数。

\exproblem{326}{52}
一个很宽的竖直水槽两相对侧壁分别有两个
截面积相同的小孔 $S=0.50\,\mathrm{cm^2}$。
两孔高度差 $\Delta h=51\,\mathrm{cm}$。
打开后两股水流向相反水平方向喷出。
求容器受到的合反冲力大小。
\exsolution
设两孔位于水面下的深度分别为 $h_1,h_2$。
由托里拆利定律，两孔出口速率满足
$v_i^2=2gh_i$。
射流质量通量为
$\dot m_i=\rho S v_i$，
容器对应的反冲力模为
$F_i=\dot m_i v_i=\rho S v_i^2
=2\rho g S h_i$。
两股喷流方向相反，反冲合力大小为
\[
F_{\rm net}=|F_1-F_2|
=2\rho gS|h_1-h_2|
=\boxed{2\rho gS\Delta h}.
\]
换算 $S=5.0\times10^{-5}\,\mathrm{m^2}$、
$\Delta h=0.51\,\mathrm m$，
\[
\boxed{F_{\rm net}
=2(1000)(9.81)(5.0\times10^{-5})(0.51)
\simeq0.50\,\mathrm N}.
\]
\excheck 水槽越深的孔喷速越大、反冲力越大，
净反冲方向指向较深孔\emph{喷射方向的反方向}。

\exproblem{327}{52}
一个宽的竖直水槽高 $H=75\,\mathrm{cm}$，
侧壁上有一条自槽底向上延伸的狭窄竖直缝，
长度为 $\ell=50\,\mathrm{cm}$、宽
$b=1.0\,\mathrm{mm}$。
开缝前水槽装满水，求刚打开缝时
全部出流水对槽壁的合反冲力大小。
\exsolution
设自槽底向上量坐标 $y\in[0,\ell]$。
此处水面下深度为 $H-y$，
出口速率由托里拆利公式为
$v(y)=\sqrt{2g(H-y)}$。
微小出流面积 $\mathrm dA=b\,\mathrm dy$。
单位时间流出的水平动量为
$\mathrm dF=\rho v(y)^2\,\mathrm dA
=2\rho g(H-y)b\,\mathrm dy$。
沿整个缝长度积分：
\[
F=\int_0^\ell2\rho gb(H-y)\,\mathrm dy
=2\rho gb(H\ell-\ell^2/2).
\]
因此
\[
\boxed{F=\rho gb\ell(2H-\ell)}.
\]
代入数值：
$\boxed{F\simeq4.91\,\mathrm N\approx5.0\,\mathrm N}$。
\excheck 这是\emph{打开瞬间}的反冲力。
随水面下降，各点深度改变，反冲力不再保持不变。

\exproblem{328}{53}
一个大水箱经内径半径
$r=0.50\,\mathrm{cm}$ 的直角弯管出水
（原书图1.87）。
弯管水平部分长 $\ell=22\,\mathrm{cm}$，
相对点 $O$ 的水平动量作用力臂为 $\ell$。
体积流量为 $Q=0.50\,\mathrm{L/s}$。
求出流水对管壁关于 $O$ 的反作用\emph{力矩大小}。
\exsolution
设管截面积 $A=\pi r^2$。
由不可压缩连续性，出水速度
$v=Q/A$。垂直向下的出口质量流率
$\dot m=\rho Q$，每秒携出的竖向动量为
$\dot m\,v=\rho Q^2/A$。
相对入口点 $O$，垂直动量通量的力臂为
$\ell$；入口水平动量通量沿通过 $O$ 的方向，
不贡献该点的力矩。
因此液体对弯管的反力矩模为
\[
\boxed{N=\ell\frac{\rho Q^2}{\pi r^2}}.
\]
取水密度 $\rho=1000\,\mathrm{kg/m^3}$、
$Q=5.0\times10^{-4}\,\mathrm{m^3/s}$、
$r=5.0\times10^{-3}\,\mathrm m$、$\ell=0.22\,\mathrm m$，
\[
\boxed{N\simeq0.70\,\mathrm{N\,m}}.
\]
\excheck 必须将升/秒换算成立方米/秒；
该式假定忽略管内流动产生的复杂附加壁面力矩与
非均匀出口速度剖面。

\exproblem{329}{53}
一个大开口水箱侧壁装有向外收缩的水平短管
（原书图1.88），入口横截面积
$S=3.0\,\mathrm{cm^2}$、出口面积
$s=1.0\,\mathrm{cm^2}$。
自由水面高出短管中心 $h=4.6\,\mathrm m$，
忽略黏性。求试图将短管从水箱侧壁
\emph{向外拔出}的水平水动力大小。
\exsolution
出水口通大气，流速
$v_s=\sqrt{2gh}$，
入口流速由连续性为
$v_S=(s/S)v_s$。
入口相对大气的静压由伯努利式为
\[
p_S-p_0=\rho gh-\frac12\rho v_S^2
=\rho gh\left(1-\frac{s^2}{S^2}\right).
\]
取收缩短管内水为控制体。
其水平动量流出减流入的通量为
$\rho s v_s(v_s-v_S)$。
入口压强、出口大气压与短管壁对流体作用力
合起来提供该动量变化。
以相对大气压计，短管\emph{对水}的水平合力为
\[
F_{{\rm wall}\to{\rm water}}
=\rho s v_s(v_s-v_S)-(p_S-p_0)S
=-\rho gh\,\frac{(S-s)^2}{S}.
\]
牛顿第三定律使水对短管向外的作用力为
\[
\boxed{F_{\rm pull}=\rho gh\frac{(S-s)^2}{S}}.
\]
代入 SI 单位 $S=3\times10^{-4}\,\mathrm{m^2}$、
$s=10^{-4}\,\mathrm{m^2}$，得
$\boxed{F_{\rm pull}\simeq6.0\,\mathrm N}$。
\excheck 不可只用静水压强乘出口面积，
还必须计入收缩管两端的动量流率变化。

\exproblem{330}{53}
一个盛水的竖直圆柱形容器以恒角速度
$\omega$ 绕自身对称轴旋转，液体最终达到
相对容器静止的稳定状态。
求：(a) 自由液面的形状；
(b) 如果槽底中心静压为 $p_0$，
槽底离轴距离 $r$ 处的静压强。
\exsolution
在随容器转动的参考系中，
液体处于相对静力平衡：
\[
\nabla p=\rho(\omega^2r\,\bm e_r-g\bm e_z).
\]
因此径向与竖直偏导数满足
$\partial_rp=\rho\omega^2r$、
$\partial_zp=-\rho g$。
积分得
\[
p(r,z)=C+\tfrac12\rho\omega^2r^2-\rho gz.
\]
自由液面压强为大气常数，
故自由表面相对其轴心的高度为
\[
\boxed{z(r)-z(0)=\frac{\omega^2r^2}{2g}},
\]
是旋转抛物面。
槽底各点同高 $z=0$，
压强分布为
\[
\boxed{p_{\rm bottom}(r)=p_0+\frac12\rho\omega^2r^2}.
\]
\excheck 自由表面上压强处处相同，
但高度随 $r$ 增加；底面同高而压强随 $r$ 增加，
二者不矛盾。

\exproblem{331}{53}
半径为 $R=10\,\mathrm{cm}$ 的水平薄圆盘
置于充满油的水平圆柱形狭缝内
（原书图1.89）。
圆盘与上、下固定平面的间隙均为
$h=1.0\,\mathrm{mm}$，
油的动力黏度
$\eta=0.08\,\mathrm{P}$
（泊，$1\,\mathrm P=0.1\,\mathrm{Pa\,s}$）。
圆盘绕中心轴以角速度
$\omega=60\,\mathrm{rad/s}$ 旋转。
忽略边缘效应，求黏滞力耗散功率。
\exsolution
盘面距轴半径为 $r$ 处线速度为
$v_\theta=\omega r$。在窄隙内，近似线性速度剖面
产生剪切应力
\[
\tau(r)=\eta\frac{\omega r}{h}.
\]
单侧盘面半径 $r$、宽 $\mathrm dr$ 的环带面积
$2\pi r\,\mathrm dr$，
其摩擦力矩微元
$\mathrm dN_1=\tau(r)\,2\pi r\,\mathrm dr\cdot r
=2\pi\eta\omega r^3\,\mathrm dr/h$。
上下两面都有黏滞作用，总力矩
\[
N=2\int_0^R\mathrm dN_1
=\frac{\pi\eta\omega R^4}{h}.
\]
所需驱动功率等于力矩乘角速度：
\[
\boxed{P=N\omega=\frac{\pi\eta\omega^2R^4}{h}}.
\]
用 $\eta=0.008\,\mathrm{Pa\,s}$ 计算，
$\boxed{P\simeq9.0\,\mathrm W}$。
\excheck 必须包含圆盘的\emph{上下两个面}，
并正确把泊换算为 SI 黏度单位。

\exproblem{332}{53}
半径为 $R_1$ 的无限长内圆柱
以恒定轴向速度 $v_0$ 在
半径为 $R_2>R_1$ 的静止同轴外圆柱内移动，
中间充满黏性液体。
忽略端部效应、假定稳定层流且无轴向压强梯度。
求液体轴向速度 $v_z(r)$ 的径向分布。
\exsolution
流动只沿轴向，速度只随半径 $r$ 变化。
牛顿黏性流体的稳定轴向运动方程在无压强梯度时为
\[
0=\eta\frac1r\frac{\mathrm d}{\mathrm dr}
\left(r\frac{\mathrm dv_z}{\mathrm dr}\right).
\]
积分一次给 $r\,\mathrm dv_z/\mathrm dr=C_1$，
再积分得到
$v_z(r)=C_1\ln r+C_2$。
利用无滑移边界
$v_z(R_1)=v_0$、$v_z(R_2)=0$，
\[
\boxed{v_z(r)=v_0
\frac{\ln(r/R_2)}{\ln(R_1/R_2)}},
\qquad R_1\le r\le R_2.
\]
\excheck 在内壁为 $v_0$、外壁为零，
中间连续单调下降；圆柱坐标中的分布是对数函数，
不是两平行平板间的严格线性函数。

\exproblem{333}{53--54}
两长同轴圆柱之间有黏度为 $\eta$ 的层流液体，
内柱半径 $R_1$ 且静止，
外柱半径 $R_2>R_1$、以角速度 $\omega_2$
绕公共轴旋转。已知半径 $r$ 处圆柱面上的剪切应力
$\tau(r)=\eta r\,\mathrm d\omega_f/\mathrm dr$，
其中 $\omega_f(r)$ 是流体局部角速度。
求：(a) $\omega_f(r)$；
(b) 单位轴向长度上作用于外柱的摩擦力矩大小。
\exsolution
稳定轴对称层流中，单位长度通过半径 $r$ 圆柱面的
扭矩为
\[
N'(r)=2\pi r^2\tau(r)
=2\pi\eta r^3\frac{\mathrm d\omega_f}{\mathrm dr}.
\]
稳态角动量传递要求 $N'(r)$ 与 $r$ 无关。
故 $\mathrm d\omega_f/\mathrm dr=C/r^3$，
积分得到
$\omega_f(r)=A+B/r^2$。
无滑移边界为
$\omega_f(R_1)=0$、
$\omega_f(R_2)=\omega_2$。
解得
\[
\boxed{\omega_f(r)=
\frac{\omega_2R_2^2}{R_2^2-R_1^2}
\left(1-\frac{R_1^2}{r^2}\right)}.
\]
由导数
$\omega_f'(r)=
2\omega_2R_1^2R_2^2/
[(R_2^2-R_1^2)r^3]$，
得到外柱受到的阻转力矩每单位长度大小
\[
\boxed{N'=
\frac{4\pi\eta\omega_2 R_1^2R_2^2}
{R_2^2-R_1^2}}.
\]
\excheck 该结果是\emph{单位轴向长度}的力矩，
量纲为 $\mathrm{N\,m/m}=\mathrm N$；
实际长为 $L$ 的外柱总阻转矩为 $N'L$。

\exproblem{334}{54}
半径为 $R$、长度为 $\ell$ 的圆管中流过
密度为 $\rho$、动力黏度为 $\eta$ 的稳定层流，
轴向流速剖面为
$v(r)=v_0(1-r^2/R^2)$。
求：(a) 体积流量；
(b) 管内流体动能；
(c) 流体对管壁的总黏滞摩擦力；
(d) 两端压强差。
\exsolution
(a) 以同心环带积分流量：
\[
Q=\int_0^R v(r)2\pi r\,\mathrm dr
=2\pi v_0\int_0^R(r-r^3/R^2)\,\mathrm dr
=\boxed{\frac12\pi v_0R^2}.
\]
(b) 管内各环带长度均为 $\ell$，动能
\[
K=\frac12\rho\ell\int_0^R v(r)^2\,2\pi r\,\mathrm dr
=\boxed{\frac16\pi\ell R^2\rho v_0^2}.
\]
(c) 牛顿黏性应力在管壁的大小
$|\tau_w|=\eta|v'(R)|=2\eta v_0/R$。
壁面积 $2\pi R\ell$，所以
\[
\boxed{F_f=|\tau_w|(2\pi R\ell)=4\pi\eta\ell v_0}.
\]
(d) 稳态充分发展流动中，
端面压强推动力平衡壁面黏滞阻力，
$\Delta p\,(\pi R^2)=F_f$，故
\[
\boxed{\Delta p=\frac{4\eta\ell v_0}{R^2}}.
\]
\excheck 与泊肃叶定律
$Q=\pi R^4\Delta p/(8\eta\ell)$
互相一致；本题 $v_0$ 是轴心最大速率，
平均速率为 $v_0/2$。
% === END 07-hydrodynamics-325-334.tex ===

% === BEGIN 07-hydrodynamics-335-339.tex blob b25c5a8e63ecca7684b6a4a236f94bcc8e04f5ac ===
% Irodov user PDF pp.63-64, original printed pp.54-55.
% Final original hydrodynamics section 1.7 problems 335-339.

\exproblem{335}{54}
原书图1.90中，水从宽大水箱 $A$
通过一根截面均匀但有黏性阻力的水平管流出。
沿管道相隔同样距离 $\ell$ 装两根竖直测压管，
其水柱高分别为 $h_2=20\,\mathrm{cm}$、
$h_1=10\,\mathrm{cm}$；
从水箱接口到第一根测压管、从第二根测压管到出口
也各为 $\ell$。水箱自由水面高出管轴
$h_3=35\,\mathrm{cm}$。
求管内水流速率 $v$。
\exsolution
在截面不变的细管内，稳定流动速度 $v$ 不随
轴向位置变化；黏性压头损失与等长管段长度成正比。
两测压管相距 $\ell$，压头损失为
$h_2-h_1=10\,\mathrm{cm}$。
因此沿上游再外推长度 $\ell$，
管入口处的压力水头为
\[
h_{\rm inlet}=h_2+(h_2-h_1)
=2h_2-h_1=30\,\mathrm{cm}.
\]
水箱自由面到管入口还可用于把水从近静止
加速到管速 $v$ 的水头为
\[
\Delta h=h_3-h_{\rm inlet}
=h_3-2h_2+h_1=5\,\mathrm{cm}.
\]
对水箱自由面和管入口应用含上述分配的能量式：
$v^2/(2g)=\Delta h$。
于是
\[
\boxed{v=\sqrt{2g(h_3-2h_2+h_1)}
=\sqrt{2(9.81)(0.05)}
\simeq1.0\,\mathrm{m/s}}.
\]
\excheck 不能把全部 $h_3$ 当作动能水头；
图示测压管读数已显示管内存在黏性沿程损失。

\exproblem{336}{54}
圆形管道的内半径沿轴向按
$r(x)=r_0e^{-\alpha x}$ 缓慢减小，
其中 $\alpha=0.50\,\mathrm{m^{-1}}$。
流体密度、黏度不变，体积流量恒定。
求沿程相隔 $\Delta x=3.2\,\mathrm m$ 的
两个截面的雷诺数比值。
\exsolution
取管内流动的特征长度为管径
$D(x)=2r(x)$，平均速度
$\bar v(x)=Q/[\pi r(x)^2]$，
雷诺数
\[
\mathrm{Re}(x)=\frac{\rho\bar v(x)D(x)}{\eta}
=\frac{2\rho Q}{\pi\eta r(x)}.
\]
因此在均匀流量下，
$\mathrm{Re}(x)\propto1/r(x)\propto e^{\alpha x}$。
两个位置 $x$ 和 $x+\Delta x$ 的比值
\[
\boxed{\frac{\mathrm{Re}(x+\Delta x)}{\mathrm{Re}(x)}
=e^{\alpha\Delta x}=e^{0.5\times3.2}
=e^{1.6}\simeq4.95\approx5}.
\]
\excheck 管径越小，平均速度增长更快于特征长度减小，
故雷诺数增大而非减小；本式以密度、黏度和流量恒定为前提。

\exproblem{337}{54}
半径为 $r_1=1.2\,\mathrm{mm}$ 的小球
在甘油中以速度不超过
$v_1=23\,\mathrm{cm/s}$ 时处于给定的层流条件。
在水中换成半径
$r_2=5.5\,\mathrm{cm}$ 的大球，
问达到相同临界雷诺数时至少需多大速度 $v_2$？
甘油与水黏度分别为
$\eta_1=13.9\,\mathrm P$、
$\eta_2=0.011\,\mathrm P$。
\exsolution
球体绕流的特征长度为直径 $2r$，
雷诺数
$\mathrm{Re}=2\rho vr/\eta$。
按题意取两个液体在\emph{同一临界雷诺数}时比较，
\[
\frac{2\rho_1v_1r_1}{\eta_1}
=\frac{2\rho_2v_2r_2}{\eta_2}.
\]
所以
\[
\boxed{v_2=v_1
\frac{r_1\rho_1\eta_2}{r_2\rho_2\eta_1}}.
\]
代入甘油密度约
$\rho_1=1.26\times10^3\,\mathrm{kg/m^3}$、
水密度 $\rho_2=1.00\times10^3\,\mathrm{kg/m^3}$
及所给黏度，得到
\[
\boxed{v_2\simeq5.0\times10^{-6}\,\mathrm{m/s}
=5\,\mathrm{\mu m/s}}.
\]
\excheck 本题按原书把雷诺数阈值称为
“层流变湍流”，但围绕球的 Stokes 流失效
与充分发展湍流并不必然是同一个转捩点；
计算结果严格属于\emph{雷诺数相似}下的阈值换算。

\exproblem{338}{54}
一颗铅球在动力黏度为
$\eta=13.9\,\mathrm P$ 的甘油中
以终端速率稳定下沉。
定义球直径为雷诺数特征长度，
当临界值 $\mathrm{Re}_{\rm cr}=0.5$ 时
视为超出题定层流范围。
求球仍符合这一范围的最大直径 $d_{\max}$。
\exsolution
记铅密度 $\rho_s$、甘油密度 $\rho_f$。
在 Stokes 阻力近似中，
球重力减浮力与黏滞阻力平衡：
\[
(\rho_s-\rho_f)\frac{\pi d^3}{6}g
=3\pi\eta d\,v_t.
\]
于是
\[
v_t=\frac{(\rho_s-\rho_f)gd^2}{18\eta}.
\]
雷诺数定义为
$\mathrm{Re}=\rho_fv_td/\eta$，
代入后
\[
\mathrm{Re}=
\frac{\rho_f(\rho_s-\rho_f)gd^3}{18\eta^2}.
\]
在 $\mathrm{Re}=\mathrm{Re}_{\rm cr}$ 处，
\[
\boxed{d_{\max}=
\left[\frac{18\mathrm{Re}_{\rm cr}\eta^2}
{\rho_f(\rho_s-\rho_f)g}\right]^{1/3}}.
\]
以 $\rho_s\simeq11.3\times10^3\,\mathrm{kg/m^3}$、
$\rho_f\simeq1.26\times10^3\,\mathrm{kg/m^3}$、
$\eta=1.39\,\mathrm{Pa\,s}$，
\[
\boxed{d_{\max}\simeq5.2\,\mathrm{mm}\approx5\,\mathrm{mm}}.
\]
\excheck 直径因公式中 $d^3$ 出现而为\emph{立方根}；
若误取平方根会破坏量纲。
雷诺数恰为0.5时 Stokes 近似属于临界估算。

\exproblem{339}{54--55}
直径 $d=3.0\,\mathrm{mm}$ 的钢球
由静止开始在橄榄油中下沉，
油的动力黏度为 $\eta=0.90\,\mathrm P$。
按 Stokes 线性阻力模型，
求钢球速率距终端速率相差不到
$n=1.0\%$ 所需时间。
\exsolution
设球体积 $V=\pi d^3/6$、钢密度 $\rho_s$、
油密度 $\rho_f$。
用向下为正写牛顿运动方程：
\[
m\dot v=(\rho_s-\rho_f)Vg-3\pi\eta d\,v,
\qquad m=\rho_sV.
\]
令
$k=3\pi\eta d$、
$\tau=m/k=\rho_s d^2/(18\eta)$。
方程的终端速率为
$v_t=(\rho_s-\rho_f)Vg/k$，
初始 $v(0)=0$ 的解是
\[
v(t)=v_t(1-e^{-t/\tau}).
\]
故相对终端速率差
$(v_t-v)/v_t=e^{-t/\tau}$。
要求它等于 $\varepsilon=0.01$ 时，
\[
\boxed{t=\tau\ln\frac1\varepsilon
=\frac{\rho_s d^2}{18\eta}\ln100}.
\]
取 $\rho_s\simeq7.8\times10^3\,\mathrm{kg/m^3}$、
$\eta=0.090\,\mathrm{Pa\,s}$，
\[
\boxed{t\simeq0.20\,\mathrm s}.
\]
\excheck 模型未计入附加质量和历史阻力。
对题给球径及橄榄油黏度，实际终端雷诺数可能不够低，
因此 Stokes 线性阻力给出的是\emph{题定模型的结果}，
并非无需校准的高精度真实沉降时间。
% === END 07-hydrodynamics-335-339.tex ===
